\documentclass[11pt]{article}

\usepackage[a4paper,margin=2.5cm]{geometry}
\usepackage{amsmath,amssymb}
\usepackage{hyperref}

\title{GRA Ramanujan Generator:\\
Multilevel Foam Functional for Symbolic Formula Discovery}

\author{oleg bitsoev}
\date{March 2026}

\begin{document}
\maketitle

\begin{abstract}
We present a small experimental playground, the \emph{GRA Ramanujan Generator}, which uses a multilevel GRA foam functional to guide evolutionary discovery of symbolic formulas for mathematical constants such as $\pi$, $e$ and $\ln 2$. The key idea is to go beyond pure numerical accuracy and simple complexity penalties, and introduce a structural ``Ramanujan-style'' prior that prefers interpretable expressions built from meaningful constants and operations.
\end{abstract}

\section{Introduction}

Symbolic regression and automated formula discovery are increasingly used in mathematics and the natural sciences to search for interpretable laws and conjectures rather than purely black-box models. Recent work such as the Ramanujan Machine and AI Feynman demonstrates that symbolic formulas can be discovered automatically and then inspected by humans.

In this note we describe a lightweight experimental system, the \emph{GRA Ramanujan Generator}, built on top of the GRA Meta-Nullification framework. The system exposes a web playground where users can watch an evolutionary process search for formulas approximating a target constant, while a multilevel GRA foam functional balances three objectives: numerical accuracy, simplicity, and structural meaning.

\section{GRA Foam Functional}

Let a candidate formula $f$ induce a numerical approximation $\hat{c} = f()$ of a target constant $c$.
We define three levels of the GRA foam functional:
\begin{itemize}
  \item Level 0 (pure error)
  \[
    \Phi^{(0)}(f) = | \hat{c} - c |.
  \]
  \item Level 1 (error + complexity)
  \[
    \Phi^{(1)}(f) = \Phi^{(0)}(f) + \lambda \cdot \mathrm{complexity}(f),
  \]
  where $\lambda > 0$ controls the penalty for expression size / tree depth.
  \item Level 2 (structural prior)
  \[
    \Phi^{(2)}(f) = \mathrm{structure\_penalty}(f),
  \]
  which rewards Ramanujan-style structure and penalises noisy constants.
\end{itemize}

Without the structural term, a standard two-level functional
\[
  \Phi_{\mathrm{GRA}}^{(2\text{-level})}(f)
  = \alpha \Phi^{(0)}(f) + (1-\alpha)\Phi^{(1)}(f),
\]
where $\alpha \in [0,1]$, tends to favour short numeric hacks: very small error and low complexity, but no mathematical meaning (e.g.\ expressions like $\sqrt{9.901}$).

To mitigate this, we introduce a three-level foam functional
\[
  \Phi_{\mathrm{GRA}}(f)
  = \alpha_0 \Phi^{(0)}(f)
  + \alpha_1 \Phi^{(1)}(f)
  + \alpha_2 \Phi^{(2)}(f),
\]
with $\alpha_0,\alpha_1,\alpha_2 \ge 0$ and $\alpha_0+\alpha_1+\alpha_2 \approx 1$.

\subsection{Structural Ramanujan prior}

The structural term $\Phi^{(2)}$ operates on the expression tree of $f$.
The current prototype uses the following heuristic components:
\begin{itemize}
  \item Rewards the presence of known constants such as $\pi$, $e$, $\ln 2$ and small rationals $p/q$ with $|p|,|q|$ small.
  \item Penalises arbitrary floating-point constants with many digits or without a simple rational approximation.
  \item Rewards certain operation patterns (powers, logs, products) that frequently appear in Ramanujan-style identities.
\end{itemize}
The overall $\Phi^{(2)}$ is a weighted sum of these features, with negative values indicating more desirable structure.

\section{Experiment: Approximating $e$}

In the playground, we run an evolutionary search for formulas approximating Euler's number $e$.
We compare two settings:

\begin{enumerate}
  \item \textbf{2-level GRA (no structural prior).}
  Here $\Phi_{\mathrm{GRA}} = \alpha \Phi^{(0)} + (1-\alpha)\Phi^{(1)}$.
  The system often converges to trivial short expressions that numerically fit the target but have no obvious meaning.

  \item \textbf{3-level GRA with Ramanujan prior.}
  Here $\Phi_{\mathrm{GRA}} = \alpha_0 \Phi^{(0)} + \alpha_1 \Phi^{(1)} + \alpha_2 \Phi^{(2)}$ with, for example,
  $\alpha_0 = 0.3$, $\alpha_1 = 0.3$, $\alpha_2 = 0.4$.
\end{enumerate}

In a typical run targeting $e$, the 3-level GRA setup produces a best formula of the form
\[
  f(x) = (4.227)^{\ln 2},
\]
with numerical value
\[
  f() \approx 2.7158753,\quad
  e \approx 2.7182818,
\]
so the absolute error is about $2.4 \cdot 10^{-3}$.

The complexity of this expression (in our tree-based metric) is $3$, and the structural term $\Phi^{(2)}$ is strongly negative because the formula uses the meaningful constant $\ln 2$ and a simple power structure, rather than arbitrary numeric constants.

Although this specific expression is not claimed to be a deep new identity, the experiment clearly shows how the structural prior can steer the search away from degenerate numeric hacks towards interpretable expressions built from familiar constants.

\section{Discussion and Outlook}

The GRA Ramanujan Generator is a small but instructive example of how a multilevel foam functional can be used to control symbolic regression:
\begin{itemize}
  \item Level~0 enforces numerical agreement with a target.
  \item Level~1 keeps expressions reasonably simple.
  \item Level~2 expresses domain knowledge about which kinds of formulas humans tend to find meaningful.
\end{itemize}

While the current playground focuses on mathematical constants, the same architecture can be applied to scientific modelling:
discovering empirical laws, soft sensors and surrogate models where interpretability matters.
Future work may integrate more principled structural priors, benchmarks and connections to existing systems such as the Ramanujan Machine and AI Feynman.

\section*{Availability}

A live demo of the GRA Ramanujan Generator is available at:
\begin{center}
\url{https://gra-ramanujan-evolut-rvxf.bolt.host/}
\end{center}
The full engine currently resides in a private research repository; this note is intended as an accompanying description of the experimental playground.

\end{document}
